% Source: https://artofproblemsolving.com/community/c3h2540106_basic_mathematics
% Stable identity: aops-topic-2540106-q01
\textbf{(a)} Divide $2x^2+3y^2=kxy$ by $y^2\ne0$ and put $t=\dfrac xy$: $2t^2+3=kt$, i.e. $2t^2-kt+3=0$. This equation does not depend on $x$ and $y$. Since it has the real root $t=\dfrac xy$, its discriminant is non-negative: $k^2-24\ge0$.

\textbf{(b)} By (a), $\dfrac xy$ always equals one of the at most two roots $t_1\le t_2$ of $2t^2-kt+3=0$. If $t_1=t_2$, the ratio is constant. If $t_1<t_2$, a continuously varying quantity that took both values would also have to pass through the values strictly between them, which are not roots; so it takes only one of them. Either way $\dfrac xy=c$ for a fixed constant $c$, i.e. $x=cy$ and $x\propto y$. $\blacksquare$

(Without some such assumption the conclusion can fail: for $k=5$ the roots are $1$ and $\dfrac32$, and the pairs $(x,y)=(1,1)$ and $(3,2)$ both satisfy $2x^2+3y^2=5xy$ with different ratios.)
